\documentclass[11pt,a4paper]{article}

\usepackage[utf8]{inputenc}
\usepackage[T1]{fontenc}
\usepackage[french]{babel}
\usepackage{lmodern}
\usepackage{geometry}
\usepackage{microtype}
\usepackage{xcolor}
\usepackage{amsmath}
\usepackage{booktabs}
\usepackage{array}
\usepackage{tabularx}
\usepackage{longtable}
\usepackage{enumitem}
\usepackage{listings}
\usepackage{fancyhdr}
\usepackage{hyperref}

\geometry{margin=2.3cm}
\hypersetup{
    colorlinks=true,
    linkcolor=blue!45!black,
    urlcolor=blue!55!black,
    pdftitle={TP n°2 -- Chiffrement hybride et RSA}
}

\definecolor{codebg}{RGB}{246,248,250}
\definecolor{codegreen}{RGB}{35,120,70}
\definecolor{codeblue}{RGB}{30,75,150}
\definecolor{codered}{RGB}{160,45,45}

\lstdefinelanguage{RustSimple}{
    morekeywords={fn,let,mut,const,pub,struct,impl,use,mod,match,if,else,return,as,true,false},
    sensitive=true,
    morecomment=[l]{//},
    morestring=[b]"
}

\lstset{
    language=RustSimple,
    backgroundcolor=\color{codebg},
    basicstyle=\ttfamily\small,
    keywordstyle=\color{codeblue}\bfseries,
    commentstyle=\color{codegreen},
    stringstyle=\color{codered},
    numbers=left,
    numberstyle=\tiny\color{gray},
    stepnumber=1,
    numbersep=8pt,
    frame=single,
    rulecolor=\color{black!15},
    breaklines=true,
    showstringspaces=false,
    columns=fullflexible,
    keepspaces=true
}

\pagestyle{fancy}
\fancyhf{}
\lhead{ECUE332 -- Protocoles de sécurité}
\rhead{Chiffrement hybride et RSA}
\cfoot{\thepage}
\setlength{\headheight}{14pt}

\setlist[itemize]{itemsep=2pt,topsep=4pt}
\setlist[enumerate]{itemsep=3pt,topsep=4pt}
\newcommand{\fichier}[1]{\texttt{#1}}
\newcommand{\commande}[1]{\texttt{#1}}
\newcommand{\reponse}{\textbf{Réponse.} }

\begin{document}
\pagenumbering{roman}

\begin{titlepage}
    \centering
    \vspace*{2cm}
    {\Large Université Paris-Est Créteil\par}
    \vspace{0.5cm}
    {\large ECUE332 -- Concepts et protocoles pour la sécurité des communications\par}
    \vspace{2.2cm}
    {\Huge\bfseries TP n°2\par}
    \vspace{0.4cm}
    {\LARGE\bfseries Chiffrement hybride et RSA\par}
    \vspace{1.5cm}
    {\large Implémentation en Rust\par}
    \vfill
    \begin{tabular}{rl}
        Étudiant : & \rule{8cm}{0.4pt} \\
        Enseignant : & J. Diaz \\
        Date : & \today
    \end{tabular}
    \vfill
\end{titlepage}

\tableofcontents
\newpage
\pagenumbering{arabic}

\section{Introduction}

L'objectif de ce TP est de comprendre le rôle des chiffrements asymétrique et
symétrique dans l'envoi d'un document confidentiel. Donald souhaite transmettre
un fichier à Vladimir sur un réseau contrôlé par un attaquant capable d'intercepter,
de modifier, de remplacer ou de rejouer les messages.

L'énoncé fournit des exemples en Python. Le travail a été réalisé en Rust avec des
bibliothèques cryptographiques existantes. Les mécanismes demandés restent les mêmes :

\begin{itemize}
    \item RSA simplifié pour comprendre les calculs mathématiques ;
    \item RSA-OAEP pour chiffrer une petite donnée de manière sûre et aléatoire ;
    \item AES-256-GCM pour chiffrer efficacement le contenu d'un fichier ;
    \item RSA-PSS pour authentifier Donald dans le mode avancé ;
    \item un identifiant unique signé pour détecter le rejeu.
\end{itemize}

Le but final est d'obtenir les propriétés suivantes : confidentialité du document,
détection des modifications, authentification de l'expéditeur et refus d'une enveloppe
déjà reçue.

\section{Organisation du projet et correspondance avec le TP}

Chaque partie du sujet correspond à un exécutable Rust distinct. Cela permet de faire
une démonstration progressive sans mélanger toutes les notions dès le début.

\begin{longtable}{p{3.6cm}p{4.1cm}p{7.1cm}}
\toprule
\textbf{Fichier} & \textbf{Partie du TP} & \textbf{Rôle} \\
\midrule
\endhead
\fichier{rsa\_simple.rs} & Rappels RSA & Calcule $n$, $\varphi(n)$ et $d$, puis chiffre et déchiffre l'entier 65. \\
\fichier{rsa\_oaep.rs} & RSA-OAEP & Génère une paire RSA-2048, chiffre avec OAEP/SHA-256 et réalise les tests demandés. \\
\fichier{hybride.rs} & RSA + AES-GCM & Chiffre \fichier{message.txt} avec AES-256-GCM et protège la clé AES avec RSA-OAEP. \\
\fichier{commun.rs} & Mode avancé & Définit l'enveloppe JSON, les AAD, le mot de passe de démonstration et les fonctions Base64. \\
\fichier{generer\_cles.rs} & Gestion des clés & Crée les paires RSA-3072 de Donald et Vladimir et les enregistre au format PEM. \\
\fichier{donald.rs} & Émetteur avancé & Chiffre le fichier, construit l'enveloppe JSON et la signe avec RSA-PSS. \\
\fichier{vladimir.rs} & Récepteur avancé & Vérifie le rejeu et la signature, déchiffre la clé AES, puis le document. \\
\fichier{main.rs} & Aide & Affiche simplement la liste des commandes disponibles. \\
\bottomrule
\end{longtable}

Les principales dépendances sont \texttt{rsa}, \texttt{aes-gcm}, \texttt{sha2},
\texttt{serde-json}, \texttt{base64} et \texttt{rand}. Elles évitent de réimplémenter
des primitives cryptographiques complexes et sensibles aux erreurs.

\section{Partie 1 -- RSA avec de petits nombres}

\subsection{Calculs réalisés}

L'énoncé impose :
\[
    p=61, \qquad q=53, \qquad e=17.
\]
Le programme calcule :
\[
    n=pq=61\times 53=3233,
\]
\[
    \varphi(n)=(p-1)(q-1)=60\times 52=3120.
\]

La clé privée contient $d$, l'inverse modulaire de $e$ modulo $\varphi(n)$ :
\[
    ed\equiv 1\pmod{\varphi(n)}.
\]
L'algorithme d'Euclide étendu donne :
\[
    d=2753, \qquad 17\times 2753\equiv 1\pmod{3120}.
\]

Pour $m=65$, le chiffrement et le déchiffrement sont :
\[
    c=m^e\bmod n=65^{17}\bmod 3233=2790,
\]
\[
    m=c^d\bmod n=2790^{2753}\bmod 3233=65.
\]

Le fichier \fichier{rsa\_simple.rs} utilise une exponentiation modulaire rapide afin
de ne jamais calculer directement les très grandes puissances. Il utilise aussi
l'algorithme d'Euclide étendu pour trouver $d$.

\begin{verbatim}
Clé publique : (17, 3233)
Clé privée : (2753, 3233)
Message initial : 65
Message chiffré : 2790
Message déchiffré : 65
\end{verbatim}

\subsection{Réponses aux questions}

\paragraph{Q1 -- Que se passe-t-il si le même message est chiffré deux fois ?}
\reponse Le même message produit toujours le même chiffré. Le RSA brut est
déterministe : avec les mêmes valeurs de $m$, $e$ et $n$, le calcul donne toujours
le même résultat.

\paragraph{Q2 -- Pourquoi cela pose-t-il un problème ?}
\reponse Si l'attaquant peut deviner les messages possibles, il peut chiffrer chaque
hypothèse avec la clé publique et comparer le résultat au chiffré intercepté. Il peut
donc reconnaître un message sans connaître la clé privée. C'est une attaque par
dictionnaire.

\paragraph{Q3 -- Pourquoi $p=61$ et $q=53$ ne protègent-ils pas des données réelles ?}
\reponse Le module $n=3233$ est beaucoup trop petit et peut être factorisé
immédiatement. En retrouvant $p$ et $q$, un attaquant calcule $\varphi(n)$ puis $d$.
Cette partie sert uniquement à comprendre les mathématiques de RSA.

\section{Partie 2 -- Chiffrement RSA-OAEP}

\subsection{Principe et implémentation}

Le fichier \fichier{rsa\_oaep.rs} répond à la partie RSA-OAEP du TP. Une paire RSA
de 2048 bits est générée pour Vladimir. La clé publique chiffre et la clé privée
déchiffre. OAEP utilise SHA-256 et ajoute notamment une graine aléatoire au message
avant l'opération RSA.

Le chiffrement est donc probabiliste : le même message peut produire plusieurs
chiffrés différents. Cette propriété corrige le principal défaut observé avec le RSA
brut. La source d'aléa \texttt{OsRng} provient du système d'exploitation.

\begin{lstlisting}[caption={Opérations principales de la partie RSA-OAEP}]
let chiffre = cle_publique.encrypt(
    &mut rng,
    Oaep::new::<Sha256>(),
    message,
)?;

let dechiffre = cle_privee.decrypt(
    Oaep::new::<Sha256>(),
    &chiffre,
)?;
\end{lstlisting}

\subsection{Réponses aux questions}

\paragraph{Q1 -- Chiffrement et déchiffrement.}
\reponse Le message est chiffré avec la clé publique de Vladimir et déchiffré avec sa
clé privée. Le programme retrouve exactement le texte initial :
\texttt{Hello Vladimir let's play GTA VI}.

\paragraph{Q2 -- Deux chiffrements du même message.}
\reponse Les deux résultats sont différents, bien que le message et la clé publique
soient identiques. OAEP ajoute un nouvel aléa à chaque chiffrement. Contrairement au
RSA brut, un attaquant ne peut donc plus reconnaître directement un message en
comparant des chiffrés.

\paragraph{Q3 -- Déchiffrement avec une autre clé privée.}
\reponse Le déchiffrement échoue. Seule la clé privée correspondant à la clé publique
utilisée au chiffrement peut retrouver le message.

\paragraph{Q4 -- Message de 200 octets.}
\reponse Le chiffrement est refusé. Avec RSA-2048 et OAEP/SHA-256, la taille maximale
du message vaut :
\[
    256-2\times 32-2=190\text{ octets}.
\]
Un message de 200 octets dépasse donc la capacité du bloc RSA.

\paragraph{Q5 -- Peut-on chiffrer directement plusieurs mégaoctets ?}
\reponse Non. RSA est limité à de petites données et ses calculs sont coûteux. Il est
adapté à la protection d'une petite clé symétrique, pas au chiffrement direct d'un
fichier. Cette limite conduit au chiffrement hybride.

\begin{verbatim}
Taille du message chiffré : 256 octets
Deux chiffrements du même message sont différents : true
Déchiffrement avec une autre clé refusé : true
Message de 200 octets refusé : true
\end{verbatim}

\section{Partie 3 -- Chiffrement hybride RSA + AES-GCM}

\subsection{Déroulement du protocole}

Le fichier \fichier{hybride.rs} réalise les quatre opérations manquantes du sujet.
Le fonctionnement est le suivant :

\begin{enumerate}
    \item Vladimir possède une paire de clés RSA-3072.
    \item Donald lit le contenu de \fichier{message.txt}.
    \item Donald génère une clé AES-256 aléatoire et un nonce de 12 octets.
    \item AES-GCM chiffre le fichier et produit un tag d'authentification.
    \item RSA-OAEP chiffre uniquement la clé AES avec la clé publique de Vladimir.
    \item Vladimir retrouve la clé AES avec sa clé privée RSA.
    \item Vladimir vérifie le tag AES-GCM et déchiffre le document.
\end{enumerate}

Les AAD (\emph{Additional Authenticated Data}) utilisées ici sont :
\begin{center}
    \path{version=1;sender=Donald;recipient=Vladimir}
\end{center}
Elles ne sont pas chiffrées,
mais elles participent au calcul du tag. Toute modification des AAD provoque donc
l'échec du déchiffrement.

\subsection{Données secrètes et données transmises}

\begin{table}[h]
\centering
\begin{tabularx}{\textwidth}{p{4.5cm}X}
\toprule
\textbf{Doit rester secret} & \textbf{Peut être transmis} \\
\midrule
Clé privée RSA de Vladimir & Clé publique RSA de Vladimir \\
Clé AES en clair & Clé AES chiffrée par RSA-OAEP \\
Document en clair & Document chiffré accompagné de son tag \\
 & Nonce AES-GCM et AAD \\
\bottomrule
\end{tabularx}
\end{table}

Le nonce n'est pas un mot de passe. Il peut être public, mais il ne doit jamais être
réutilisé avec la même clé AES-GCM.

\subsection{Réponses aux questions}

\paragraph{Q1 -- Les quatre opérations complétées.}
\reponse Le programme chiffre le contenu avec AES-GCM, chiffre la clé AES avec
RSA-OAEP, déchiffre la clé AES avec la clé privée de Vladimir, puis déchiffre le
contenu avec AES-GCM.

\paragraph{Q2 -- Vérification du fichier retrouvé.}
\reponse La comparaison entre le tableau d'octets initial et le tableau retrouvé renvoie
\texttt{true}. Le programme écrit le résultat dans le fichier
\path{message_dechiffre.txt}.

\paragraph{Q3 -- Quelles données doivent rester secrètes ?}
\reponse La clé privée RSA, la clé AES en clair et le document en clair doivent rester
secrets. La clé publique, le nonce, les AAD, le contenu chiffré avec son tag et la clé
AES chiffrée peuvent circuler sur le réseau.

\paragraph{Q4 -- Modification d'un octet.}
\reponse Le programme inverse un bit du premier octet chiffré. Le déchiffrement est
refusé, car le tag recalculé par AES-GCM ne correspond plus au tag reçu. Cette
expérience montre que GCM apporte l'intégrité en plus de la confidentialité.

\paragraph{Q5 -- Pourquoi le nonce ne doit-il pas être réutilisé ?}
\reponse Avec une même clé et un même nonce, GCM réutilise le même flot de
chiffrement. Un attaquant peut alors obtenir une relation entre les messages en
calculant $C_1\oplus C_2=M_1\oplus M_2$. La sécurité de l'authentification peut
également être compromise.

\paragraph{Q6 -- Pourquoi AES pour le fichier et RSA pour la clé ?}
\reponse AES est rapide et accepte des données de taille quelconque. RSA est beaucoup
plus lent et limité par la taille de son module, mais il permet de protéger une clé sans
secret partagé préalable. L'association des deux apporte efficacité et distribution sûre
de la clé de session.

\section{Synthèse de la partie obligatoire}

\subsection{Clé utilisée pour chaque opération}

\begin{table}[h]
\centering
\begin{tabularx}{\textwidth}{X p{7cm}}
\toprule
\textbf{Opération} & \textbf{Clé utilisée} \\
\midrule
Chiffrer la clé AES à destination de Vladimir & Clé publique RSA de Vladimir \\
Retrouver la clé AES reçue par Vladimir & Clé privée RSA de Vladimir \\
Chiffrer le fichier avec AES-GCM & Clé AES de session \\
Déchiffrer le fichier avec AES-GCM & Même clé AES de session \\
\bottomrule
\end{tabularx}
\end{table}

\paragraph{Q2 -- L'attaquant peut-il retrouver la clé AES ?}
\reponse Non, dans les conditions normales de sécurité. La clé publique permet de
chiffrer mais pas d'inverser RSA-OAEP. Sans la clé privée de Vladimir, l'attaquant ne
peut pas retrouver la clé AES.

\paragraph{Q3 -- AES-GCM prouve-t-il que Donald est l'expéditeur ?}
\reponse Non. Un tiers connaît aussi la clé publique de Vladimir. Il peut créer sa
propre clé AES, chiffrer un document valide avec AES-GCM et protéger cette clé avec la
clé publique de Vladimir. Le tag prouve seulement que l'enveloppe n'a pas été modifiée
depuis sa création par une personne connaissant la clé AES.

\paragraph{Q4 -- Comment vérifier l'origine du message ?}
\reponse Donald doit signer les éléments transmis avec sa clé privée de signature.
Vladimir vérifie ensuite cette signature avec une copie authentifiée de la clé publique
de Donald. Le mode avancé utilise RSA-PSS pour cela.

\section{Mode avancé}

\subsection{A -- Enveloppe JSON et Base64}

\paragraph{Q1 -- Structure de l'enveloppe.}
\reponse La structure définie dans \fichier{commun.rs} contient les champs suivants :

\begin{lstlisting}[language={},caption={Structure logique de l'enveloppe JSON}]
{
  "version": 1,
  "id": "identifiant unique en Base64",
  "expediteur": "Donald",
  "destinataire": "Vladimir",
  "cle_aes_chiffree": "...Base64...",
  "nonce": "...Base64...",
  "contenu_chiffre": "...Base64...",
  "signature": "...Base64..."
}
\end{lstlisting}

Le champ \texttt{contenu\_chiffre} contient le texte chiffré suivi du tag produit par
la bibliothèque AES-GCM.

\paragraph{Q2 -- Pourquoi utiliser Base64 ?}
\reponse JSON représente naturellement du texte, mais les clés, nonces, chiffrés et
signatures sont des octets arbitraires. Base64 les convertit en caractères facilement
stockables dans JSON. Base64 n'apporte aucune confidentialité : toute personne peut le
décoder.

\paragraph{Q3 -- Fonctions de conversion.}
\reponse \fichier{commun.rs} définit \texttt{vers\_base64} pour encoder les octets et
\texttt{depuis\_base64} pour les retrouver. Les données encodées par Donald sont
décodées par Vladimir avant les opérations cryptographiques, ce qui vérifie le trajet
aller-retour.

\paragraph{Q4 -- Modification après décodage Base64.}
\reponse Dans le protocole avancé, la signature est vérifiée avant le déchiffrement :
une modification du champ chiffré provoque donc d'abord le message
\og signature invalide \fg{}. Si l'on isole le test AES-GCM, ou si l'enveloppe est
resignée, la vérification du tag échoue également. Base64 ne protège pas la donnée ;
ce sont la signature et le tag qui détectent la modification.

\subsection{B -- Séparation des programmes et gestion des clés}

\paragraph{Q1 -- Clés nécessaires à chaque programme.}

\begin{table}[h]
\centering
\begin{tabularx}{\textwidth}{p{2.7cm}X X}
\toprule
\textbf{Programme} & \textbf{Clé privée} & \textbf{Clé publique} \\
\midrule
Donald & Clé privée de Donald pour signer & Clé publique de Vladimir pour chiffrer la clé AES \\
Vladimir & Clé privée de Vladimir pour déchiffrer la clé AES & Clé publique authentifiée de Donald pour vérifier la signature \\
\bottomrule
\end{tabularx}
\end{table}

\paragraph{Q2 -- Protection PEM par phrase secrète.}
\reponse \fichier{generer\_cles.rs} enregistre les clés privées au format PKCS\#8 PEM
chiffré. Le mot de passe du TP est écrit dans le code uniquement pour simplifier la
démonstration. Le programme essaie aussi un mauvais mot de passe et constate que le
chargement est refusé. Dans une vraie application, le mot de passe ne serait pas codé
en dur.

\paragraph{Q3 -- Pourquoi conserver les clés ?}
\reponse Une paire RSA représente une identité et doit rester stable pour pouvoir
déchiffrer ou vérifier des données créées précédemment. Si Vladimir remplace sa paire
avant de traiter une ancienne enveloppe, sa nouvelle clé privée ne correspond plus à la
clé publique utilisée par Donald : l'ancienne clé AES devient indéchiffrable.

\paragraph{Q4 -- Permissions Linux adaptées.}
\reponse Le dossier de clés privées peut être limité avec \commande{chmod 700 cles}
et chaque clé privée avec \commande{chmod 600 fichier\_prive.pem}. Seul le
propriétaire peut alors lire ou modifier la clé. Une clé publique peut être en
\commande{644}, car elle n'est pas secrète.

\paragraph{Q5 -- Substitution de la clé publique de Vladimir.}
\reponse Donald chiffrerait la clé AES avec la clé publique de l'attaquant. Celui-ci
pourrait alors retrouver le document, tandis que Vladimir ne pourrait pas le déchiffrer.
Il faut authentifier la clé publique, par exemple avec un certificat, une PKI, une
empreinte vérifiée par un autre canal ou un magasin de clés de confiance.

\subsection{C -- Signature RSA-PSS}

\paragraph{Q1 -- Données signées.}
\reponse Donald signe la version, l'identifiant, l'expéditeur, le destinataire, la clé
AES chiffrée, le nonce et le contenu chiffré avec son tag. Ainsi, un attaquant ne peut
pas remplacer séparément un élément ou une métadonnée sans invalider la signature.

\paragraph{Q2 -- Pourquoi exclure la signature elle-même ?}
\reponse Sa valeur n'existe pas encore au moment de calculer les données à signer.
L'inclure créerait une dépendance circulaire : il faudrait connaître la signature pour
calculer cette même signature.

\paragraph{Q3 -- Sérialisation déterministe.}
\reponse Donald et Vladimir reconstruisent une structure Rust dédiée,
\texttt{DonneesSignees}, dont l'ordre des champs est fixe. Elle est sérialisée en JSON
de la même manière des deux côtés. Les espaces et l'ordre du fichier JSON externe ne
modifient donc pas les octets reconstruits pour la vérification.

\paragraph{Q4 -- Mauvaise clé ou champ modifié.}
\reponse Une autre clé publique ne correspond pas à la clé privée ayant créé la
signature : la vérification échoue. De même, modifier un champ signé change son
empreinte SHA-256 et provoque le rejet. Ces échecs indiquent que l'identité attendue ou
l'intégrité des champs n'est pas correcte.

\paragraph{Q5 -- Pourquoi ne pas fournir la clé publique dans le même JSON ?}
\reponse Un attaquant pourrait créer sa propre paire, modifier le message, le signer
avec sa clé privée et placer sa clé publique dans le JSON. La signature serait
mathématiquement valide, mais elle ne prouverait pas l'identité de Donald. La clé
publique de Donald doit être obtenue par une source authentifiée et indépendante.

\subsection{D -- Rejeu et analyse de sécurité}

\paragraph{Q1 -- Le tag ou la signature détectent-ils un rejeu ?}
\reponse Non. Une ancienne enveloppe valide conserve une signature et un tag valides.
Ces mécanismes détectent une modification, pas une nouvelle transmission des mêmes
octets.

\paragraph{Q2 -- Mécanisme anti-rejeu.}
\reponse Donald génère un identifiant aléatoire de 16 octets pour chaque enveloppe.
Cet identifiant est inclus dans la signature. Vladimir conserve les identifiants déjà
acceptés dans \fichier{identifiants\_acceptes.txt}. Une enveloppe dont l'identifiant est
déjà présent est refusée. L'identifiant n'est mémorisé qu'après le succès de toutes les
vérifications.

\paragraph{Q3 -- Pourquoi une date signée ne suffit-elle pas ?}
\reponse Une enveloppe peut être rejouée plusieurs fois pendant sa période de validité.
De plus, les horloges peuvent être décalées. La date aide à limiter une fenêtre
d'acceptation, mais il faut aussi mémoriser un identifiant unique, un nonce de protocole
ou un compteur.

\paragraph{Q4 -- Compromission ultérieure de la clé privée de Vladimir.}
\reponse Oui, l'attaquant peut déchiffrer les anciens échanges enregistrés. Il utilise la
clé privée volée pour retrouver chaque ancienne clé AES protégée par RSA-OAEP, puis
utilise cette clé AES avec le nonce pour déchiffrer le document. Ce protocole ne fournit
pas de confidentialité persistante, appelée aussi \emph{forward secrecy}.

\paragraph{Q5 -- Pourquoi séparer clés de signature et de déchiffrement ?}
\reponse La séparation limite les conséquences d'une compromission, permet des
durées de vie et des rotations différentes, évite les confusions d'usage et facilite
l'application de politiques précises. En production, une clé ne devrait servir qu'à un
usage clairement défini.

\section{Résultats et démonstration}

\subsection{Commandes}

Le code actuel lit \fichier{message.txt} avec le chemin relatif
\texttt{../message.txt}. Les commandes suivantes sont donc lancées depuis le dossier
\fichier{src} :

\begin{verbatim}
cd src

cargo run --bin rsa_simple
cargo run --bin rsa_oaep
cargo run --bin hybride

cargo run --bin generer_cles
cargo run --bin donald
cargo run --bin vladimir
\end{verbatim}

Pour montrer l'anti-rejeu, il suffit de lancer une seconde fois :

\begin{verbatim}
cargo run --bin vladimir
\end{verbatim}

La deuxième réception de la même enveloppe est refusée, car son identifiant figure déjà
dans \fichier{identifiants\_acceptes.txt}. Pour créer une nouvelle enveloppe avec une
nouvelle clé AES, un nouveau nonce et un nouvel identifiant, il faut relancer
\commande{cargo run --bin donald}.

\subsection{Résultats significatifs}

\begin{itemize}
    \item RSA simplifié retrouve le message 65 après avoir produit le chiffré 2790.
    \item Deux chiffrements OAEP du même texte sont différents.
    \item Une mauvaise clé privée et un message OAEP de 200 octets sont refusés.
    \item Le fichier déchiffré par le mode hybride est identique à l'original.
    \item La modification d'un seul bit du contenu AES-GCM est détectée.
    \item Le mode avancé vérifie la signature de Donald avant de déchiffrer.
    \item Une seconde réception de la même enveloppe est rejetée comme rejeu.
\end{itemize}

\section{Limites de la démonstration}

Ce projet est pédagogique et ne doit pas être considéré comme une application prête
pour la production. Le mot de passe des clés privées est codé en dur, les identifiants
anti-rejeu sont stockés dans un simple fichier texte, et la confiance dans les clés
publiques est supposée. Une application réelle utiliserait un gestionnaire de secrets,
un stockage transactionnel, des certificats ou une PKI, des clés distinctes selon leur
usage et un protocole fournissant éventuellement la confidentialité persistante.

\section{Conclusion}

Le TP montre pourquoi RSA ne doit pas servir à chiffrer directement un gros fichier.
RSA-OAEP protège efficacement une petite clé de session, tandis qu'AES-256-GCM chiffre
rapidement le document et détecte ses modifications. Cette combinaison constitue le
chiffrement hybride.

Le mode avancé complète la confidentialité et l'intégrité par une signature RSA-PSS,
qui permet à Vladimir de vérifier l'origine de l'enveloppe à condition de posséder une
clé publique authentifiée de Donald. Enfin, l'identifiant unique signé et mémorisé
permet de détecter le rejeu, ce que ni le tag AES-GCM ni la signature ne peuvent faire
seuls.

\end{document}
